\documentclass[10pt,a4paper]{amsart}
\usepackage[margin=19mm]{geometry}
\usepackage{amsmath,amssymb,mathtools}
\usepackage[scr=rsfs,cal=euler]{mathalfa}
\usepackage{xfrac}
\usepackage[unicode,hidelinks,bookmarks=false]{hyperref}
\newcommand{\ie}{i.e.\ }
\newcommand{\cf}{cf.\ }
\newcommand{\resp}{respectively\ }
\newcommand{\Subsection}{\subsection}
\providecommand{\nobreakdash}{\nobreak\hbox{-}}
\setlength{\emergencystretch}{2em}
\multlinegap0pt
\usepackage{bm}
\usepackage{tikz}
\usepackage{amssymb}
\usepackage{xcolor}
\usepackage[normalem]{ulem}
\usepackage{mathtools}
\newtheorem{lemma}{Lemma}
\newtheorem{proposition}{Proposition}
\title{von Neumann Inequality for a class of Doubly Contractive Weighted Shifts}
\author{Soumyadip Dey, Rajeev Gupta,, Surjit Kumar}
\date{Source: arXiv:2604.06978v2 (CC BY 4.0)}
\begin{document}
\maketitle

\noindent\emph{Source and licence.} von Neumann Inequality for a class of Doubly Contractive Weighted Shifts. Version arXiv:2604.06978v2 (CC BY 4.0), distributed by arXiv under the Creative Commons Attribution 4.0 International licence, http://creativecommons.org/licenses/by/4.0/, as recorded in the arXiv licence metadata for this exact version and shown on the abstract page https://arxiv.org/abs/2604.06978v2. Full licence notice: \texttt{licenses/arxiv-2604.06978-CC-BY-4.0.txt}.\par\smallskip

\noindent\textit{Editorial scope.} Counted unit: the article's lemma shift (source0.tex lines 560--563). The article's complete original proof of it is delivered with it (source0.tex lines 582--613, one \texttt{proof} environment of 32 lines, which the article places away from the statement). No mathematics is written, repaired or re-derived: the statement and the proof are the article's own lines, in the article's own notation, and no retained line is substituted or re-flowed.  Retained uncounted as the source-local context the proof needs: lines 465--467; lines 506--512; lines 530--533.  Nothing was omitted from any retained span, so the package carries no omission record.  No retained line contains a citation, so the wrapper carries no bibliography and no bibliography entry.  No retained line contains a figure environment, an image-inclusion command or drawing code, so no image is delivered and the package declares no asset dependency.  Editorial formatting only, in addition: an AMS \texttt{amsart} preamble that restates the article's own declarations and macros used by the retained lines, namely the environments \texttt{lemma}, \texttt{proposition} on separate counters, together with the standard packages those lines need.  The retained environments print in this document as Lemma 1, Proposition 1 in order of appearance, whereas the article prints the same items with the numbering of its own full document.  The complete native source of the article as distributed in the arXiv e-print for this exact version is bundled unchanged as \texttt{sources/198059/source0.tex}.
\begin{lemma}\label{Shilov}
The Shilov boundary of \( X \) is contained in $X_0.$
\end{lemma}

\begin{proposition}\label{Arveson}
    Let \(X \subset\mathbb{C}^N  \) be compact, and suppose that \(T: X\to \mathcal{B}(\mathcal{H)}^d\) is an analytic function such that \(T(z)\) is a d-tuple of commuting bounded operators  for all \(z\in X\). Then the following statements are true: 
    \begin{enumerate}
        \item If the tuple \(T(z)\) satisfies  von Neumann's inequality over $\mathbb B_d$ for all \(z\in \partial_0 X\) then $T(z)$ satisfies von Neumann's inequality over $\mathbb B_d$ for all $z\in  X$.
        \item If the tuple \(T(z)\) dilates to a tuple of commuting spherical unitaries for all \(z\in \partial_0 X\) then $T(z)$ dilates to a tuple of commuting spherical unitaries for all $z\in X$.
    \end{enumerate}
\end{proposition}

\begin{proposition}\label{Dilation for X_0}
    Let $\boldsymbol T=(T_1, \ldots, T_d)$ be a doubly contractive $d$-variable weighted shift with weights $\{w_{I,j}: I\in\mathbb{Z}_+^d, j\in\{1,\ldots, d\}\}$ satisfying $w_{I,j}=m_j$ for some real numbers $m_1,\ldots,m_d,$ and $\sum_{j=1}^{d} m_j^2 \leq 1.$
Then $\boldsymbol T$ admits a spherical dilation. In particular, $\boldsymbol T$ satisfies von Neumann’s inequality.
\end{proposition}

\begin{lemma}\label{Sperical unitary dilation for truncated weighted shift}
Suppose $m_1,\ldots, m_d$ are positive real numbers such that \(\sum_{j=1}^d m_j^2 \leq1\).
Let $\boldsymbol T=(T_1,\ldots, T_d)$ be a \( d \)-variable doubly contractive truncated weighted shift with weights $(w_{I,j})$ such that \( w_{I,j}\in (0, m_j] \) for all $I\in\mathbb{Z}_+^d$ and $j\in\{1,\ldots, d\}.$ Then $\boldsymbol T$  dilates to a \( d \)-tuple of commuting spherical unitaries.
\end{lemma}

\begin{proof}[Proof of Lemma \ref{Sperical unitary dilation for truncated weighted shift}:]
Define \[
r := \left( 
\max \left\{ 
\frac{|w_{I,j}|}{m_j} :  (I, j) \text{ is scalable} 
\right\}
\right)^{-1}.
\]
Then \( r> 1\).  Let \( D_r(0) \subset \mathbb{C}\) be the  closed unit disc of radius \(r\)  around 0 . For \(t\in D_r(0)\) and \( 
(I,j)\in \mathcal{I} \) , define 
\[
\hat{w}_{I,j}(t) = 
\begin{cases}
t w_{I,j} & \text{if } (I,j) \text{ is scalable }, \\
w_{I,j} & \text{otherwise},
\end{cases}
\]
and let \( \hat{w}(t) = (\hat{w}_{I,j}(t))_{(I,j) \in I} \). Then it follows from the maximum modulus principle that there exists \( t_0 \in \partial D_r(0) \) with
\[
|f(w)| = |f(\hat{w}(1))| \leq |f(\hat{w}(t_0))|,\] so setting \( \hat{w} = \hat{w}(t_0) \), iterating this process finitely many times yields a point \( v \in {X_0} \) such that all  weights are  \(m_j\) in respective direction  such that \( |f(w)| \leq |f(v)| \) and also satisfying commutativity without changing row and column contraction. So we need to show that \( \hat{w}(t) \) is a commuting family, that is
\[
\hat{w}_{I,j}(t) \hat{w}_{I+\varepsilon_j,k}(t) = \hat{w}_{I,k}(t) \hat{w}_{I+\varepsilon_k,j}(t)
\]
for all \( t \in D_r(0) \) and all multi-indices \( I \) with \( |I| \leq N-1 \) and \( 1 \leq j, k \leq d \).
Let \( I \) be such a multi-index. If \( I \) is bad, it follows from (a) that \( I + \varepsilon_j \) and \( I + \varepsilon_k \) are bad as well, and hence no pairs in \( I \) which appear in the above equation are scalable. If \( I \) and \( I + \varepsilon_j + \varepsilon_k \) are good, then it follows again from (a) that no pairs in the equation are scalable. Thus, it remains to consider the case where \( I \) is good and \( I + \varepsilon_j + \varepsilon_k \) is bad. In this case, exactly one of \( (I,j) \) and \( (I + \varepsilon_j, k) \) is scalable, depending on whether \( I + \varepsilon_j \) is good or bad. Similarly, exactly one of \( (I,k) \) and \( (I + \varepsilon_k, j) \) is scalable. Thus,
\[\hat{w}_{I,j}(t) \hat{w}_{I+\varepsilon_j,k}(t) = t w_{I,j} w_{I+\varepsilon_j,k} = t w_{I,k} w_{I+\varepsilon_k,j} = \hat{w}_{I,k}(t) \hat{w}_{I+\varepsilon_k,j}(t),
\]
as stated. 
Now the proof of Lemma \ref{Shilov} is complete.

Note that Proposition \ref{Dilation for X_0} shows that any weighted shift in $X_0$ has spherical dilation. In particular, using Lemma \ref{Shilov}, we have spherical dilation for any weighted shift in $\partial_0 X.$ An application of Proposition \ref{Arveson} now completes the proof of lemma.
\end{proof}

\end{document}
